%%%%%%%%%%%%%%%%%%%%%%%%%%%%%%%%%%%%%%%%%
% 中文简历 - Chinese Resume
% Based on Medium Length Professional CV Template
%%%%%%%%%%%%%%%%%%%%%%%%%%%%%%%%%%%%%%%%%

%----------------------------------------------------------------------------------------
%	文档配置
%----------------------------------------------------------------------------------------

\documentclass{resume} % 使用自定义 resume.cls 样式
\usepackage{ctex} % 中文支持
\usepackage{hyperref}
\hypersetup{colorlinks=true,linkcolor=blue,filecolor=blue,urlcolor=blue,citecolor=cyan}
\usepackage[left=0.75in,top=0.6in,right=0.75in,bottom=0.6in]{geometry} % 页面边距
\newcommand{\tab}[1]{\hspace{.2667\textwidth}\rlap{#1}}
\newcommand{\itab}[1]{\hspace{0em}\rlap{#1}}
\name{王一通} % 姓名

\address{+44 07900492737\\ +86 13265139183 \\ yitong.7.wang@kcl.ac.uk \\ ytwangcs@gmail.com} % 联系方式
\address{研究方向：机器人学习、人机交互} % 研究兴趣
\begin{document}

%----------------------------------------------------------------------------------------
%	教育背景
%----------------------------------------------------------------------------------------

\begin{rSection}{教育背景}
{\bf 香港科技大学，中国香港特别行政区} \hfill {\em 2026年9月 - 至今} 
\\ 机械工程博士，导师：{\href{https://yanggao.people.ust.hk/index.html}{高阳教授}}
\\ 隶属{\href{https://yanggao.people.ust.hk/CAIRSS.html}{太空可持续性人工智能机器人中心}}

{\bf 伦敦国王学院，英国伦敦} \hfill {\em 2023年9月 - 2026年1月} 
\\ 医疗技术研究型硕士（MRes）
\\ 导师：{\href{https://profiles.ucl.ac.uk/60656-erika-molteni/about}{Erika Molteni博士}}、{\href{https://www.kcl.ac.uk/people/michela-antonelli}{Michela Antonelli博士}} 和 {\href{https://www.kcl.ac.uk/people/thomas-booth}{Thomas Booth博士}} 
\\ GPA：72/100（预计：Distinction）

{\bf 华南理工大学，中国广州} \hfill {\em 2018年9月 - 2022年6月} 
\\ 工学学士（导师：{\href{https://www2.scut.edu.cn/ft_en/2025/1021/c45325a606216/page.htm}{张东教授}}） \hfill GPA：3.15/4 或 82.48/100
\\ 主修：过程装备与控制工程 \hfill 辅修：计算机科学与技术

\end{rSection}

%----------------------------------------------------------------------------------------
%	交流经历
%----------------------------------------------------------------------------------------

\begin{rSection}{交流经历}

{\bf 新加坡国立大学计算机学院暑期工作坊 2022} \hfill {\em 2022年5月 - 2022年8月}\\
方向：深度学习与嵌入式系统（导师：{\href{https://www.comp.nus.edu.sg/cs/people/ctank/}{Colin博士}}） \hfill 成绩：A

{\bf 帝国理工学院大师课"机器人、人工智能与物联网"} \hfill{\em 2021年1月 - 2021年2月}
\\ 主题：消防巡检机器人（导师：{\href{https://www.imperial.ac.uk/people/benny.lo}{Penny Lo博士}}）[{\em 最佳项目}] \hfill 成绩：95/100

\end{rSection}

%----------------------------------------------------------------------------------------
%	获奖情况
%----------------------------------------------------------------------------------------

\begin{rSection}{获奖情况} \itemsep -5pt
\item 2021年 IEEE国际机器人与自动化大会（{\bf\em ICRA}）人工智能挑战赛 {\bf 三等奖}
\item 2021年 RoboMaster机甲大师赛（中国大学生机器人大赛）{\bf 一等奖}
\item 国家级大学生创新创业训练计划（优秀结题，负责人）
\item 2020年 RoboCon四足机器人仿真竞赛 {\bf 一等奖}
\item 华南理工大学优秀学生暨奖学金（前25\%）
\item 第32届全国中学生物理竞赛（CPhO）{\bf 二等奖}
\item 第31届全国中学生物理竞赛（CPhO）{\bf 三等奖}

\end{rSection}

%----------------------------------------------------------------------------------------
%	论文发表
%----------------------------------------------------------------------------------------

\begin{rSection}{论文发表} \itemsep -5pt
{\bf 在投论文}
\item Yitong Wang, Aaliyah Adesida, Jiaheng Wang, Shiyang Lu, Ben Jackson, Thomas Booth, 《急性脑卒中远程操作机器人取栓工具包》，2025年。[\href{https://docs.google.com/document/d/1d-0TzQVWwxyS-562f4fhBHqsg136TUnLxxGucI2UqpA/edit#heading=h.8kry2pqa1z3y}{早期报告}]

\end{rSection}

%----------------------------------------------------------------------------------------
%	项目经历
%----------------------------------------------------------------------------------------

\begin{rSection}{项目经历}
{\bf 基于多模态深度学习的心理压力与孤独感检测}
\\ 负责开发能够接受多模态信号输入的Transformer深度学习模型，进一步探索不同信号模态缺失对输出结果的影响，以及提升模型鲁棒性的方法。

{\bf 急性脑卒中远程操作机器人取栓工具包[{\em 相关文章预计2025年发表}]}
\\ 负责基于SOFA框架使用PPO实现机械取栓的自动化导航。(1) 在SOFA框架中使用PPO强化学习开发机械取栓自动化导航系统，并将学习模型迁移到真实机器人。(2) 设计算法与人工操作性能对比实验。(3) 开发用于人机实验的机器人交互控制器。

{\bf 智能植物培养箱}
\\ 负责深度学习模块及软硬件接口开发。工作包括：收集植物病虫害和气象信息数据集，使用Inception V3进行植物病虫害检测，使用LSTM评估生长环境（温度、湿度、光照）。为工作坊展示制作海报和视频。

{\bf 基于规则的双机器人强化学习任务协作}\\
{\em ICRA 2021人工智能挑战赛项目（三等奖）}\\
1){\bf 定位}：基于YOLO V4算法实现机器人本体识别（Mosse滤波跟踪增强鲁棒性和稳定性），通过PCL坐标变换实现定位，准确率高于98.84\%。2){\bf 决策}：使用A3C和行为树实现机器人的事件决策（攻击、补给、移动和防御等）。3){\bf 导航}：选择激光雷达作为传感器，基于图优化使用Cartographer和AMCL算法实现机器人自身定位，分别使用A*和TEB算法实现3Hz和40Hz频率的全局和局部规划。4){\bf 云台自动瞄准}：使用OpenCV进行颜色阈值分割和分离，卡尔曼滤波预测对手运动（比基于学习的方法更快），识别准确率超过90\%，4米内误差不超过5cm，使用PID算法控制云台。

{\bf 大学生研究计划（SRP）项目}\\
主题：{\em 仿生四足机器人结构与算法研究}\\
{\em RoboCon 2021四足机器人仿真竞赛项目}
\\ 对仿生机器人腿部结构进行动力学分析，使用Fusion360应用计算科学方法实现仿生优化。该过程在不影响机械强度的情况下成功将结构重量减轻30\%。同时学习并重写了控制算法。

\end{rSection}

%----------------------------------------------------------------------------------------
%	工作经历
%----------------------------------------------------------------------------------------

\begin{rSection}{工作经历}

\begin{rSubsection}{华南理工大学，广州}{\em 2022年9月 - 2023年6月}{研究助理}{}
\item 负责数据处理、仿真与环境搭建、模型训练、数据可视化代码
\end{rSubsection}

\begin{rSubsection}{博泰机器人科技有限公司，佛山}{\em 2021年1月 - 2021年5月}{算法研究实习生}{}
\item 机器人创业公司算法实习。从事移动机器人感知算法研究，特别是视觉SLAM和多传感器融合。
\end{rSubsection}

\end{rSection}

%----------------------------------------------------------------------------------------
%	技能特长
%----------------------------------------------------------------------------------------

\begin{rSection}{技能特长}

\begin{tabular}{ @{} >{\bfseries}l @{\hspace{6ex}} l }
建模与分析 \ & AutoCAD, SolidWorks, Fusion360, ANSYS, 3D打印 \\
软件工具 & MS Office, LaTeX, Adobe系列软件 \\
开发框架 \ & PyTorch, TensorFlow, OpenCV, ROS, SOFA（软体机器人仿真）\\
编程语言 \ & Python, C++, HTML, JavaScript, MATLAB\\
语言能力 \ & 中文（母语）, 英文
\end{tabular}

\end{rSection}

\end{document}
